\section{Обсуждение}
\label{sec:Discussion}

\subsection{Сравнение предложенных подходов}

В работе рассмотрена задача восстановления зависимости между $X$ и $Y$ по группам, внутри
которых исходное попарное соответствие утрачено. Для её решения разработаны два подхода,
различающиеся параметризацией скрытого распределения, моделью условной плотности и способом
обучения.

В нейросетевом подходе CatVAE условное распределение задаётся декодером, а вариационное
распределение перестановки --- энкодером с Gumbel--Sinkhorn-преобразованием. Теоретическая
NELBO выводится непосредственно из точного группового правдоподобия с одной перестановкой
$\sigma\in S_M$ и не использует факторизацию по локальным назначениям
(Раздел~\ref{ssec:nelbo}). В реализации непрерывная Gumbel--Sinkhorn-релаксация и
построчный KL образуют дифференцируемый суррогат, позволяющий совместно обучать энкодер и
декодер методом SGD. После обучения новая группа обрабатывается прямым проходом энкодера и
детерминированной Sinkhorn-проекцией без повторной оптимизации параметров.

Во втором подходе условный закон ограничивается глобальной смесью сдвигов, а
ответственности назначения и компоненты вычисляются явно. Обучение выполняется
модифицированным EM (Раздел~\ref{ssec:em}): компонентная часть E-шага имеет аналитический
вид, распределение перестановки приближается Sinkhorn-планом, а веса и периодические сдвиги
обновляются явными формулами на M-шаге. Нейронные сети, SGD и обратное распространение в
этом подходе не используются.

Различие результатов двух подходов нельзя сводить только к различию оптимизаторов. CatVAE
использует более гибкий класс условных распределений и амортизирует вывод, но требует
обучения нейронных сетей. Модифицированный EM использует сильное предположение о глобальной
смеси сдвигов и благодаря этому получает интерпретируемые аналитические обновления. Его
преимущество на синтетическом бенчмарке отражает в том числе совпадение параметризации с
порождающей моделью.

\subsection{Интерпретация экспериментальных результатов}

Эксперименты показывают, что предложенные подходы восстанавливают не только плотность, но и
скрытую структуру сопоставления. При $M=2$ CatVAE достигает точности $98.0\%$ при
$\epsilon=0.01$ по сравнению со случайным уровнем $50\%$; с ростом шума точность снижается
до $89.7\%$ при $\epsilon=0.05$ (Рис.~\ref{fig:recovery} и
Таблица~\ref{tab:recovery-acc}). Восстановленные пары располагаются вдоль двух мод
порождающего распределения, что подтверждает способность модели распутывать перестановку,
а не только аппроксимировать маргинальные распределения наблюдаемых множеств.

По $L_2$-ошибке плотности CatVAE на рассмотренном бенчмарке остаётся одного порядка с
энтропийным транспортным оператором, но не демонстрирует устойчивого превосходства над ним
(Рис.~\ref{fig:varying-M}). Нейросетевой подход платит за более общий класс декодеров и
амортизацию вывода. Модифицированный EM при малом и умеренном шуме и малых или средних $M$
заметно превосходит референс. Этот результат получен в благоприятном режиме, где
предположение о глобальной смеси сдвигов совпадает с генератором данных, и не должен
переноситься на произвольные условные распределения без дополнительной проверки.

Рост числа групп $N$ в среднем уменьшает ошибку обоих подходов. Для модифицированного EM
при $\epsilon\in\{0.01,0.025\}$ нескольких десятков групп оказывается достаточно, чтобы
оценить сдвиги и веса смеси с малой ошибкой (Рис.~\ref{fig:varying-N}). Следовательно,
групповые непарные наблюдения действительно содержат информацию о совместной зависимости:
при подходящей модели накопление независимых групп позволяет восстановить параметры, хотя
ни в одной группе правильные пары не предъявляются алгоритму.

\subsection{Границы применимости}

Основное вычислительное ограничение связано с размером группы $M$. Оба подхода избегают
явного перебора $M!$ перестановок, но используют матрицу попарных оценок из $M^2$ элементов.
В CatVAE она обрабатывается Gumbel--Sinkhorn-релаксацией, а в модифицированном EM ---
Sinkhorn-аппроксимацией E-шага. Стоимость одной итерации Sinkhorn имеет порядок $O(M^2)$,
а качество мягкой перестановки зависит от температуры и числа нормализаций. Поэтому
практическое применение ограничено группами, для которых матрица оценок помещается в память
и Sinkhorn-проекция остаётся устойчивой.

Второе фундаментальное ограничение --- идентифицируемость сопоставления. При
$\epsilon=0.05$ компоненты условной смеси сильнее перекрываются, точность восстановления пар
падает, а ошибка EM выходит на плато около $1.1$, которое не устраняется увеличением $N$.
В таком режиме несколько перестановок почти одинаково согласуются с наблюдениями, поэтому
невозможность уверенно восстановить конкретное соответствие является свойством данных, а
не только недостатком алгоритма.

Теоретическая вариационная постановка CatVAE использует точную скрытую перестановку, однако
практическая функция обучения содержит Gumbel--Sinkhorn-релаксацию, одну выборку для
реконструкционного члена и построчный KL-суррогат. В модифицированном EM точный
перестановочный E-шаг заменён энтропийным транспортным планом, а круговое моментное
обновление сдвига не обязано максимизировать точное правдоподобие обёрнутой нормальной
смеси. Поэтому для него не действует стандартная гарантия монотонного роста
правдоподобия классического EM. Результаты также зависят от инициализации, температуры и
числа итераций Sinkhorn.

Экспериментальная проверка ограничена синтетическими одномерными данными на торе с
известными $K$ и $\epsilon$. Она не устанавливает качество подходов на многомерных или
реальных данных.

\subsection{Открытые задачи}

Для больших групп требуется уменьшить квадратичные затраты на матрицу назначений, например
за счёт разреженных кандидатных соответствий или иерархического сопоставления, не разрушая
условие биективности. Другой открытый вопрос --- построение более точной и при этом
вычислимой аппроксимации апостериорного распределения по перестановкам как для
Gumbel--Sinkhorn-обучения, так и для EM E-шага.

Для модели плотности естественным продолжением является совместное оценивание дисперсии,
числа компонент и формы компонент вместо фиксации $K$ и $\epsilon$. Важна также проверка
способности амортизированного энкодера переноситься между разными размерами групп и уровнями
шума. Наконец, переход к многомерным пространствам, другим геометриям и реальным данным
позволит установить, сохраняются ли наблюдаемые преимущества при отклонении от модели
глобальных сдвигов. Отдельная теоретическая задача состоит в формулировке условий
идентифицируемости совместного распределения по непарным группам.

\subsection{Основные результаты}

Основные результаты, полученные в работе, состоят в следующем.
\begin{enumerate}
  \item Выведена вариационная нижняя граница непосредственно для точной скрытой
  перестановки группы, без факторизации по локальным назначениям.
  \item Разработан нейросетевой CatVAE: энкодер с Gumbel--Sinkhorn-релаксацией и декодер
  условной обёрнутой гауссовой смеси совместно обучаются стохастическим градиентным методом
  и обеспечивают амортизированный вывод для новых групп.
  \item Разработан отдельный подход с явной параметризацией глобальной смеси сдвигов,
  обучаемый модифицированным EM с Sinkhorn-аппроксимацией E-шага и аналитическими
  обновлениями весов и периодических сдвигов.
  \item Экспериментально исследовано влияние шума, размера групп, объёма данных и
  Sinkhorn-температуры на качество обоих подходов.
\end{enumerate}
